\section{Hecke eigenpoints, their full determinants, and the original ordinary output}
\label{sec:hecke-eigenpoints-v44}

This chapter reconstructs the full original \texttt{TCoreEig} theorem
and the entire \texttt{OrdCoreH.Out} record.  An arbitrary arithmetic
point of the specified Hecke algebra gives one continuous semisimple
representation, its complete determinant law, and the prescribed point
property.  At a good point it also gives the actual ordered ordinary
flags and potential crystallinity at every original local bridge.
The algebra, character labels, coefficient maps, and level are unchanged.

Fix the original data $L,\Omega,\ell,n,c,H$ and $H.\mathrm{Ax}$.
Write $O=c.O$, $\overline E=\overline{\mathbb Q}_\ell$, and fix the
given algebraic isomorphism $\iota:\overline E\simeq\mathbb C$.
Let $A$ have the original class $\mathrm{IsCO}$, let
$H.\mathrm{IsWt}(\theta)$ hold, and fix the residual operator labels
$a$ and the original level $K$.  Put
\[
 N=H.MK(\theta,a,K),\qquad T=H.TK(\theta,a,K)
       \subset\operatorname{End}_A(N).
 \tag{HEP.1}
\]
We retain the source assumptions that the specified operators commute
on $N$ and that $T$ is of class $\mathrm{IsCO}$ over $O$.
For an original operator index $\alpha$, write $t_\alpha$ for its
actual generator in $T$.  Suppose given
\[
 D\in\mathrm{GDet}_n(T,H.G_q),\qquad
 \psi_{v,i}:D_v\longrightarrow T^\times,
 \qquad H.\mathrm{IsUniv}(\theta,a,K,\mathfrak m_T,D,\psi).
 \tag{HEP.2}
\]
Here $D_v=\operatorname{Dec}(H.\mathrm{pl}(v))$, and all local maps
are the source maps.  No domain assumption or flatness of a point
map is imposed on $A$ or $T$.

The original $\mathrm{LiteratureRankNHypothesis}$ is retained.
The specific arithmetic foundations used below are its global Artin
reciprocity and local idele injectivity clauses, with the local--global
and norm normalizations of class field theory; its local rank-one,
unlinked principal-series, and Langlands laws; and the original
$\mathrm{RecTExists}$.  These are also the explicit foundations of
the earlier CIP, AFS, GUA, and OFL packets.  OFL16 carries precisely
its stated period foundations.  None is relabeled as a newly proved
class field theory, local Langlands, or period comparison theorem here.

\subsection{Continuity and equality of complete determinant laws}

% BEGIN HEP1
\paragraph{HEP1. Every original algebra point preserves adic continuity.}
\label{claim:HEP1}
Let $R$ be any $O$-algebra of class $\mathrm{IsCO}$ and
$y:R\to_O\overline E$ any algebra point.  If
$f:\Gamma\to R$ is adically continuous for $\mathfrak m_R$,
then $y\circ f$ is continuous for the original topology of
$\overline E$.  In particular every evaluated determinant
coefficient and every evaluated $\psi_{v,i}$ in (HEP.2) is continuous.
The evaluation of $H.\mathrm{epsV}(R,v)$ is the original
cyclotomic character $\varepsilon_\ell|_{D_v}$.
\begin{proof}
ASP2, applied to the very point $y$, gives $0<\eta<1$ with
\[
             |y(b)|_\ell\leq\eta^r
                  \qquad(b\in\mathfrak m_R^r).
 \tag{HEP.3}
\]
Fix $g\in\Gamma$ and $\epsilon>0$.  Choose $r$ with
$\eta^r<\epsilon$.  Adic continuity says that, on a neighborhood
of $g$, $f(g')-f(g)\in\mathfrak m_R^r$.  Thus
$|y(f(g'))-y(f(g))|_\ell<\epsilon$ there.  This proves continuity
at every $g$ without a continuity assumption on the point itself.

The \texttt{cont} and \texttt{\(\psi\)\_cont} fields of
$\mathrm{IsUniv}$ supply these adic hypotheses.  The resulting
character into $\overline E^\times$ is continuous in the topology
on units: its value is continuous by the argument just given, and
its inverse value is the value at $g^{-1}$, hence is continuous too.
Finally $\mathrm{epsV}$ is defined by the structural maps
$\mathbb Z_\ell\to O\to R$ applied to $\mathrm{cycZ}$.
Since $y$ is an $O$-algebra map, evaluation is the same structural
map to $\overline E$.  This proves the stated cyclotomic identity.
\end{proof}
% END HEP1

% BEGIN HEP2
\paragraph{HEP2. Dense characteristic polynomials determine the whole evaluated law.}
\label{claim:HEP2}
Let $\Gamma$ be a topological group and let $D_1,D_2$ be
degree-$n$ determinants over a characteristic-zero field $E$,
with its Hausdorff topology.  Suppose their characteristic
coefficients on every group element are continuous.  If their
characteristic polynomials agree on a dense subset of $\Gamma$,
then $D_1=D_2$ as polynomial laws, after every commutative
$E$-algebra extension.
\begin{proof}
Each difference of corresponding coefficient functions is continuous.
Its zero set is closed and contains a dense set, so it is all of
$\Gamma$.  In particular the two trace functions agree on every
group element.

We give the algebra argument that upgrades this assertion to full
polynomial laws.  For any multiplicative determinant $\Delta$,
write
\[
 \Delta(1-tz)=\sum_{k=0}^n(-1)^k\lambda_k(z)t^k,
 \qquad \lambda_0=1,\quad \operatorname{tr}_\Delta=\lambda_1.
 \tag{HEP.4}
\]
These expressions are natural in every scalar algebra.  In a ring
of dual numbers,
$(1+\epsilon x)(1+\epsilon z)=1+\epsilon(x+z)$.
Multiplicativity therefore makes the coefficient of $\epsilon$
additive in $x,z$; substitution $\epsilon\mapsto b\epsilon$
makes it linear over the scalar algebra.  Thus
$\operatorname{tr}_\Delta$ is a natural linear map.  Agreement on
$\Gamma$ implies agreement on every element of $B[\Gamma]$, for
every $E$-algebra $B$, by finite linear expansion in group elements.

For completeness, multiplicativity also gives the Newton identities
in the form needed here.  In formal power series, or in their
arbitrary finite truncations, factor
\[
 1-(t+\epsilon)z=(1-tz)
       \left(1-\epsilon z(1-tz)^{-1}\right).
\]
Apply the natural law to this identity and compare the coefficient
of $\epsilon$.  Since $(1-tz)^{-1}=\sum_{r\geq0}t^rz^r$, it gives
\[
 \frac{d}{dt}\Delta(1-tz)
   =-\Delta(1-tz)\sum_{r\geq0}
                        \operatorname{tr}_\Delta(z^{r+1})t^r.
\]
Each coefficient is justified in a finite truncation, so no
convergence or continuity of a polynomial law is involved.
Equating coefficients yields
\[
 k\lambda_k(z)=\sum_{j=1}^k(-1)^{j-1}
                   \lambda_{k-j}(z)\operatorname{tr}_\Delta(z^j).
 \tag{HEP.5}
\]
Every positive integer is invertible in $B$, because $B$ is an
$E$-algebra.  Induction on $k$ now identifies all the $\lambda_k$
of $D_1$ and $D_2$ on every $B[\Gamma]$.  Homogeneity identifies
$\lambda_n(z)$ with the value of the degree-$n$ law on $z$.
Hence the laws are equal on every scalar algebra, as claimed.
For $n=0$ this last assertion is the constant unital law and is
immediate.  In its application here, division by integers occurs
only after passage to $\overline E$; no integer is inverted in
$O,A$, or $T$.
\end{proof}
% END HEP2

% BEGIN HEP3
\paragraph{HEP3. A spherical eigenvector identifies the full specialized determinant.}
\label{claim:HEP3}
Let $y:T\to_O\overline E$, let $\lambda,\xi$ be any original
weight and coefficient representation, and let
$Y:H.\mathrm{PartAltAlt}(\iota,\lambda,\xi)$.
Suppose $0\ne F\in Y.Y$ satisfies
\[
 H.\operatorname{opC}(\iota,\lambda,d,(w,j))F
       =\iota(y(t_{(w,j)}))F
       \quad(w\in H.W_p,\ 0\leq j<n).
 \tag{HEP.6}
\]
Then, as full determinant laws,
\[
                    D\otimes_{T,y}\overline E
                             =\det(Y.r).
 \tag{HEP.7}
\]
No arithmetic or good-point assumption is required here.
\begin{proof}
The original index $(w,j)$ denotes the spherical operator of
degree $j+1$.  Its action on $Y.Y$ is also the scalar
$Y.a(w,j+1)$ by the part's \texttt{sph} field.
Since $F\ne0$, equality of the two scalar multiples implies
\[
       \iota\bigl(y(H.\operatorname{genW}(w,k))\bigr)=Y.a(w,k)
                                      \quad(1\leq k\leq n).
\]
At $k=0$ both sides equal one by the literal definition of
$\mathrm{genW}$ and the part's \texttt{a\_zero} field.

Scalar extension of a determinant sends its characteristic
polynomial to the coefficientwise image of that polynomial:
apply naturality to the group-algebra element $X\cdot1-g$.
The same is immediate for the finite sum defining
$\mathrm{heckePol}$.  Thus the \texttt{frob} fields of
$\mathrm{IsUniv}$ and $Y$ give equality of the characteristic
polynomials in (HEP.7) at every original element
$H.\mathrm{prj}(\sigma^{-1})$, for every allowed spherical place,
prolongation, and arithmetic Frobenius $\sigma$.
The fixed $\iota$ is injective, so equality after its coefficient
map gives equality over $\overline E$ itself.

The actual set of these inverse Frobenius elements is dense in
$H.G_q$ by the original $H.\mathrm{Ax.dense}$ field.
The evaluated coefficients of $D$ are continuous by HEP1.
The coefficients of the characteristic polynomial of $Y.r(g)$
are continuous finite polynomials in its matrix entries, because
$Y.\mathrm{cont}$ holds.  HEP2 therefore proves (HEP.7), with
its full scalar-algebra quantifier.
\end{proof}
% END HEP3

\subsection{Recovering each ordered local character}

At $v\in H.V$ put $k=L_{H.S_l(v)}$, $q=\#\kappa_k$, and
$b=H.\mathrm{art}(v):W_v\to k^\times$.
For any integer-valued function $h$ on the global embeddings put
\[
 A_h(z)=\prod_{\tau:k\hookrightarrow\overline E,\ \tau\ \text{continuous}}
                       \tau(z)^{h(\tau|_L)}.
 \tag{HEP.8}
\]
This is the literal $\mathrm{algPow}$ on nonzero arguments.
Write $l_i=\operatorname{lowW}(\lambda,i)$ and
$h_i=\operatorname{htW}(\lambda,i)=l_i+i$.

% BEGIN HEP4
\paragraph{HEP4. Algebraic factors and the full Weil cyclotomic formula.}
\label{claim:HEP4}
The products (HEP.8) are finite, nonzero, and multiplicative, with
$A_{h+h'}=A_hA_{h'}$ and $A_{-h}=A_h^{-1}$.
For every original $s\in W_v$,
\[
 \varepsilon_\ell(s)=q^{\deg(s)}A_1(b(s)),\qquad
 A_{-h_i}(z)A_{l_i}(z)A_1(z)^i=1.
 \tag{HEP.9}
\]
Both statements use the given original Artin map and embedding set.
\begin{proof}
WRP1 identifies the continuous embeddings with the embeddings of
the finite extension $k/\mathbb Q_\ell$.  There are finitely many,
and every field image of $z\ne0$ is nonzero.  The elementary
integer-power rules give all assertions about $A_h$, including
the second identity in (HEP.9), since $-h_i+l_i+i=0$.
The product $A_1$ is exactly $N_{k/\mathbb Q_\ell}$.

Here is the cyclotomic normalization also for nonzero Weil degree.
Over $\mathbb Q_\ell$, arithmetic reciprocity on a unit $u$
acts on $\ell$-power roots of unity by $u^{-1}$, as verified by
the principal-idele argument in CIP5.  Arithmetic reciprocity
on the uniformizer $\ell$ acts trivially on these roots: global
reciprocity kills the positive principal idele $\ell$, its positive
real component acts trivially, and every other finite component
is a unit at a place where the cyclotomic extension is unramified.
Consequently geometric reciprocity on any
$z\in\mathbb Q_\ell^\times$ acts by
$z\ell^{-\operatorname{ord}_\ell z}$.

Local norm compatibility therefore makes geometric reciprocity on
$x\in k^\times$ act by
\[
       N_{k/\mathbb Q_\ell}(x)\,
             \ell^{-f\operatorname{ord}_k(x)}
            =A_1(x)q^{-\operatorname{ord}_k(x)},
\]
where $f=[\kappa_k:\mathbb F_\ell]$.
The original $\mathrm{IsLocalArtin}$ map is the geometric map:
its \texttt{glob} field agrees with that map under the original
global reciprocity map, and the retained local idele injectivity
makes the two values equal.  This comparison works for every
original number field and place; it does not require a split CM
place.  Equivalently it is the normalization used in the source's
\texttt{epsArt}.  Since
$\operatorname{ord}_k(b(s))=-\deg(s)$, the displayed formula is
the first identity of (HEP.9).  This retains every admissible
Artin map and extends the inertia identity of CIP5 by the required
factor $q^{\deg(s)}$.
\end{proof}
% END HEP4

% BEGIN HEP5
\paragraph{HEP5. Unit values and ordered partial products determine the local characters.}
\label{claim:HEP5}
Let $\phi_i:W_v\to\overline E^\times$ be characters and
$\omega_i:k^\times\to\mathbb C^\times$ smooth characters.
Suppose, whenever $b(s)$ is a unit,
\[
 \phi_i(s)\varepsilon_\ell(s)^i A_{l_i}(b(s))
                             =\iota^{-1}(\omega_i(b(s))),
 \tag{HEP.10}
\]
and, whenever $b(s)=\varpi$ for the source uniformizer,
\[
 \iota\left(\prod_{i\leq j}\phi_i(s)\varepsilon_\ell(s)^i\right)
  =\iota\left(\prod_{i\leq j}A_{l_i}(\varpi)\right)^{-1}
                                  \mathrm{psEig}(\omega,j+1)
                      \quad(0\leq j<n).
 \tag{HEP.11}
\]
Then for every $i$ and $s\in W_v$,
\[
       \phi_i(s)=\iota^{-1}(\omega_i(b(s)))A_{-h_i}(b(s)).
 \tag{HEP.12}
\]
\begin{proof}
If $b(s)$ is a unit, the degree is zero.  Substitute
$\varepsilon_\ell(s)=A_1(b(s))$ from HEP4 in (HEP.10),
and divide by its nonzero algebraic factors.  This gives
(HEP.12) on all such $s$.

If $b(s)=\varpi$, its degree is $-1$.  By definition,
\[
       \mathrm{psEig}(\omega,j+1)
                  =\prod_{i\leq j}\omega_i(\varpi)q^{-i}.
 \tag{HEP.13}
\]
Thus (HEP.11), pulled back through $\iota$, equates the ordered
partial products of $\phi_i(s)\varepsilon_\ell(s)^i$ and
$A_{l_i}(\varpi)^{-1}\iota^{-1}(\omega_i(\varpi))q^{-i}$.
All factors are nonzero.  The equality for $j=0$ identifies the
first factors; dividing the equality for $j$ by that for $j-1$
identifies each subsequent factor.  HEP4 says
$\varepsilon_\ell(s)=q^{-1}A_1(\varpi)$, so cancellation gives
(HEP.12) also on every preimage of the uniformizer.

The right side of (HEP.12) is a multiplicative character of $W_v$.
Choose $s_0$ with $b(s_0)=\varpi$, using the actual Artin
surjectivity.  For any $s$, let $m=\operatorname{ord}_k(b(s))$.
Then $b(ss_0^{-m})$ is a unit and
$s=(ss_0^{-m})s_0^m$.  Agreement on the first factor and on $s_0$
proves agreement on $s$.  No commutativity of $W_v$ is required.
The ordered initial products have recovered the individually
labeled characters, not merely their multiset.
\end{proof}
% END HEP5

% BEGIN HEP6
\paragraph{HEP6. The universal unit and uniformizer relations give the exact OrdChar.}
\label{claim:HEP6}
Fix $y:T\to_O\overline E$ and put $x=y|_A$.
Suppose smooth characters $\omega_i$ at $v$ satisfy
\begin{align}
 \omega_i(u)\,\iota\bigl(H.\mathrm{finPart}(\theta,x,\lambda,v,i,u)\bigr)
       &=1 &&(u\in\mathcal O_k^\times), \tag{HEP.14}\\
 \iota(y(t_{(v,j)}))
       &=\operatorname{lowVal}(\iota,\lambda,v,j+1)^{-1}
                                      \mathrm{psEig}(\omega,j+1).
                                                        \tag{HEP.15}
\end{align}
For every $i$, the literal
\[
 \mathrm{OrdChar}\bigl(\iota,b,\omega_i,h_i,
                                      y_*\psi_{v,i}\bigr)
 \tag{HEP.16}
\]
holds on the original decomposition group.
\begin{proof}
Put $t_i(u)=\operatorname{image}_T\theta(e_{v,i}(u))$ and
$\phi_i=y_*\psi_{v,i}$.  HEP1 proves continuity of $\phi_i$
and identifies the evaluated cyclotomic character.
The \texttt{\(\psi\)\_unit} field in (HEP.2) gives, for an
original Weil element mapping to $u$,
\[
                 \phi_i(s)\varepsilon_\ell(s)^i y(t_i(u))=1.
\]
The definition of the finite part and (HEP.14) give
\[
                  y(t_i(u))=A_{l_i}(u)\,
                                   \iota^{-1}(\omega_i(u))^{-1}.
\]
Combining these identities gives (HEP.10).

The \texttt{\(\psi\)\_unif} field says that the partial product
$\prod_{i\leq j}\phi_i(s)\varepsilon_\ell(s)^i$ equals
$y(t_{(v,j)})$ whenever $b(s)=\varpi$.
The source definition
\[
 \operatorname{lowVal}(\iota,\lambda,v,j+1)
            =\iota\left(\prod_{i\leq j}A_{l_i}(\varpi)\right)
\]
turns (HEP.15) into (HEP.11).  HEP5 now proves the entire
Weil-value identity of $\mathrm{OrdChar}$.
Together with smoothness of $\omega_i$ and the continuity already
proved, these are its three exact fields.  Its \texttt{OnArt}
formulation allows any decomposition-group element whose underlying
automorphism equals the given Weil element; such elements are equal
in the subgroup, so the same identity verifies that full quantifier.
\end{proof}
% END HEP6

% BEGIN HEP7
\paragraph{HEP7. The full original core eigenvalue theorem.}
\label{claim:HEP7}
Under precisely the original hypotheses of
\texttt{core\_eigen\_of\_weight\_char\_hecke\_level}, for every
$y,\lambda,\xi,Y,F$ satisfying (HEP.6) with $0\ne F\in Y.Y$,
the determinant equality (HEP.7) holds.  Moreover, at every $v$,
every smooth family $\omega_i$ satisfying (HEP.14)--(HEP.15)
gives all the original characters (HEP.16).
\begin{proof}
HEP3 gives the determinant conclusion using only the spherical
eigenvector assumptions in the source.  For an arbitrary $v$ and
an arbitrary family in the theorem's second conclusion, HEP6
gives each $\mathrm{OrdChar}$.  The original
$H.\mathrm{Ax}$ supplies $H.S_l(v)\mid\ell$, the validity and
surjectivity of $H.\mathrm{art}(v)$, and the dense inverse-Frobenius
set.  The two complete local coefficient-ring assumptions supply
HEP1 at $T$; $H.\mathrm{IsWt}$ and all remaining hypotheses are
kept exactly as in the endpoint.  This proves its entire
conjunction, without adding arithmetic, good-point, domain,
nonvanishing-module, or flatness hypotheses.
\end{proof}
% END HEP7

\subsection{The original representation at every arithmetic point}

% BEGIN HEP8
\paragraph{HEP8. Each admitted arithmetic point is realized in one original part.}
\label{claim:HEP8}
Let $P(\lambda,r)$ be any original point predicate and assume
$H.\mathrm{AutomorphicHeckeSet}(\iota,P)$.
For every admitted $\lambda$ and every $y:T\to_O\overline E$
whose restriction $x=y|_A$ is arithmetic of weight $\lambda$,
there are the original specialization data, one original part $Y$,
and a nonzero vector
\[
 F\in Y.Y\cap H.\mathrm{OrdAt}
                   (\theta,\iota,x,\lambda,\xi,K,c_0)
 \tag{HEP.17}
\]
with $P(\lambda,Y.r)$ such that every original generator acts on
$F$ by $\iota(y(t_\alpha))$.  In particular $r=Y.r$ is continuous,
semisimple, satisfies $P(\lambda,r)$, and has determinant (HEP.7).
\begin{proof}
ASP13 constructs the full original specialization datum at the
specified arithmetic $x$ and weight $\lambda$, using the unchanged
coefficient isomorphism and original operators.  AFS19(3), applied
to this datum and the given point $y$ extending $x$, places a
nonzero simultaneous eigenvector in one original part with the
property $P$.  That assertion permits the literal depth $c_0$;
equivalently one may keep its existential original depth $e\geq c_0$.
It gives (HEP.17) and the eigenvalues for both spherical and local
positive generators.  The part fields give continuity and
semisimplicity of $Y.r$.  Its spherical eigenvalues satisfy (HEP.6),
so HEP7 gives the complete determinant equality.  Goodness was
not used in choosing this representation.
\end{proof}
% END HEP8

% BEGIN HEP9
\paragraph{HEP9. The same part gives the exact de Rham polynomial at a good point.}
\label{claim:HEP9}
Suppose additionally that the point in HEP8 is good.  At every $v$
there is a smooth family $\omega_{v,i}$ satisfying
(HEP.14)--(HEP.15), pairwise distinct on local units, such that
\[
 Y.\sigma_l(v)\simeq\mathrm{PS}(\omega_v),
 \qquad
 \mathrm{DRat}\left(\iota,H.\mathrm{ovr}(v),r|_{D_v},
     \tau\mapsto\{h_i(\tau)\}_i,
     s\mapsto\prod_i(X-\omega_{v,i}(b(s)))\right).
 \tag{HEP.18}
\]
The polynomial identity has the original Artin map $b$.
\begin{proof}
The nonzero $F$ in (HEP.17) makes the part's intersection with
the corresponding ordinary space nonzero.  AFS17 applies to
that same part and supplies smooth $\omega_v$, its principal-series
equivalence, the inverse finite-part equations (HEP.14), and the
scalar action of each local positive operator on the entire
intersection.  Compare that action on $F$ with its given
generator eigenvalue.  Since $F\ne0$, the two scalars are equal,
which is exactly (HEP.15), with its unchanged $j+1$ index.

If two unit characters $\omega_{v,i},\omega_{v,j}$ were equal,
their nonzero values and (HEP.14) would make the corresponding
finite parts equal at every unit.  This contradicts the distinctness
field of $H.\mathrm{Good}$.  Thus the unit characters are pairwise
distinct.

Choose an original correspondence $\mathrm{rc}$ using
$\mathrm{RecTExists}$.  We check the local polynomial used in
the \texttt{AtLT} field.  By the original definition,
$\mathrm{PS}(\omega_v)$ is normalized induction of
$\omega_{v,i}|\cdot|^{(n-1)/2}$.  Its irreducible representation
has these singleton cuspidal support terms.  Distinct unit
restrictions exclude every linking by a unit norm twist.
The singleton Langlands calculation in GUA2, with the inverse
twist $(1-n)/2$ of $\mathrm{algNorm}$, therefore gives the actual
parameter
\[
 \mathrm{rc}(\mathrm{algNorm}(Y.\sigma_l(v)))
                   =\bigoplus_i(\omega_{v,i}\circ b_{\mathrm{rc}},0).
 \tag{HEP.19}
\]
The Artin map $b_{\mathrm{rc}}$ belonging to this correspondence
and the source $b$ both satisfy the full original
$\mathrm{IsLocalArtin}$ predicate.  Their global reciprocity
identities and the retained local idele injectivity give
$b_{\mathrm{rc}}=b$, just as in HEP4.
This verifies the needed forward principal-series calculation on
the exact source map; no arbitrary subquotient converse is invoked.

The characteristic polynomial of (HEP.19) is the product in
(HEP.18).  Every character kills commutators, so at every Weil
commutator that polynomial is $(X-1)^n$.  The rank is $n$,
hence this is exactly the additional $\mathrm{AbW}$ test in
the original $\mathrm{AtLT}$ field.  Apply $Y.\mathrm{atL}(v)$
to the chosen smooth family, its distinct unit restrictions,
the principal-series equivalence, and this allowed correspondence.
It gives $\mathrm{DRat}$ for its parameter polynomial; (HEP.19)
is precisely the replacement giving (HEP.18).  All three
$\mathrm{DRat}$ clauses, and every original inducing bridge,
remain in that assertion.
\end{proof}
% END HEP9

% BEGIN HEP10
\paragraph{HEP10. The complete original arithmetic-point field Out.pt.}
\label{claim:HEP10}
For every admitted $\lambda$ and every arithmetic point
$y:T\to_O\overline E$, the representation $r$ of HEP8 satisfies
\[
 \mathrm{Continuous}(r),\quad \mathrm{IsSSRep}(r),\quad
 P(\lambda,r),\quad D\otimes_{T,y}\overline E=\det r.
 \tag{HEP.20}
\]
If that same point is good, then for every $v$,
\[
 \mathrm{HasFlag}\bigl(r\circ H.\mathrm{dec}(v),
                                  (y_*\psi_{v,i})_i\bigr)
       \quad\text{and}\quad
 \mathrm{PotCrys}\bigl(H.\mathrm{pl}(v),r\circ H.\mathrm{dec}(v)\bigr).
 \tag{HEP.21}
\]
This is the entire literal $\mathrm{Out.pt}$ field.
\begin{proof}
HEP8 proves (HEP.20) for every arithmetic point and chooses $r$
before a good-point hypothesis is imposed.  If the point is good,
HEP9 supplies its smooth local families, distinctness, and exact
$\mathrm{DRat}$ polynomial.  HEP7 applies (HEP.14)--(HEP.15)
to supply $\mathrm{OrdChar}$ for the original evaluated characters
$y_*\psi_{v,i}$.

For each embedding $\tau$ and $i<j$, dominance of $\lambda(\tau)$
gives
\[
 \lambda(\tau)_{n-1-i}\leq\lambda(\tau)_{n-1-j}.
\]
Consequently
$h_j(\tau)-h_i(\tau)
 =\lambda(\tau)_{n-1-j}-\lambda(\tau)_{n-1-i}+j-i>0$.
The weights are strictly increasing in the precise source order.
The local representation is continuous because $r$ is continuous
and $H.\mathrm{dec}(v)$ is the original continuous local map.
The original $H.\mathrm{Ax}$ gives the odd-prime condition,
the place-over-$\ell$ condition, and the exact Artin validity.

All hypotheses of OFL16 are therefore satisfied on the unchanged
local representation with these individually labeled characters.
It gives its actual ordered flag and potential crystallinity for
every original bridge, which is (HEP.21).  In particular no
semisimplification of the local representation replaces it.
The prime $3$ is included.  The implication to (HEP.21) retains
the original good-point guard, whereas (HEP.20) holds on every
admitted arithmetic point as required.
\end{proof}
% END HEP10

% BEGIN HEP11
\paragraph{HEP11. The whole literal OrdCoreH.Out record.}
\label{claim:HEP11}
Retain every original hypothesis of
\texttt{core\_hecke\_of\_automorphic\_hecke\_dense\_weight},
including the supplied original automorphic-Hecke-set predicate
and an admitted $\lambda_0$ with $H.\mathrm{DenseWt}(\theta,\lambda_0)$.
Then the original
\[
           H.\mathrm{Out}(\theta,a,K,P,\mathfrak m_T,D,\psi)
 \tag{HEP.22}
\]
holds, with all of its \texttt{univ}, \texttt{sep}, and \texttt{pt}
fields and their original quantifiers.
\begin{proof}
The \texttt{univ} field is exactly the given (HEP.2), with no
replacement determinant or local character.
For \texttt{sep}, fix any admitted weight $\lambda$ satisfying
$\mathrm{DenseWt}$, not merely the initially supplied $\lambda_0$.
AFS20 applied to the same $T$ says that
\[
 \left[\forall y:T\to_O\overline E,\quad
 H.\mathrm{Good}(\theta,y|_A,\lambda)\Rightarrow y(t)=0\right]
                      \Longrightarrow t=0
 \tag{HEP.23}
\]
for every $t\in T$.  This is exactly the full \texttt{sep} field.
Finally HEP10 proves \texttt{pt} for every admitted weight and
every arithmetic point.  Its representation is chosen before its
conditional good-point conclusion, as the target record requires.
These three fields prove (HEP.22).  The initially supplied
$\lambda_0$ is retained as in the source theorem; no argument
restricts its later universally quantified weights to $\lambda_0$.
\end{proof}
% END HEP11

% BEGIN HEP12
\paragraph{HEP12. Installation for the original unitary Hecke data.}
\label{claim:HEP12}
Let the original $\mathrm{HypothesesH'}$, Hida setting, residual
datum $\bar r$, and weight object $\Lambda$ be given, together with
\[
                 R_H:\mathrm{UnitaryHeckeData}(Q,\iota,\Lambda).
\]
Keep any original level $K$, set
$H=hset(Q,R_H.\tau_0)$, $A=\Lambda.\Lambda$,
$\theta=R_H.\theta$, $a=R_H.a$, and let $T$ be exactly (HEP.1).
Under the original commutativity and $\mathrm{IsCO}(O,T)$ hypotheses,
every supplied (HEP.2) has the full output
\[
 H.\mathrm{Out}\bigl(\theta,a,K,
       \mathrm{UnitaryPointProperties}(\bar r,\iota),
       \mathfrak m_T,D,\psi\bigr).
 \tag{HEP.24}
\]
Thus its representation at each admitted arithmetic point carries
the original polarization, every Teichm\"uller-family inertia type,
conditional cuspidal conclusion, and full automorphic witness;
at a good point it also has exactly the prescribed flags and
potential crystallinity.  Equation (HEP.23) holds at every admitted
dense weight, including the zero weight.
\begin{proof}
FNL16 supplies every original $H.\mathrm{Ax}$ field at this
$hset$.  The original data record supplies
$R_H.\mathrm{isWt}$ and the literal $\mathrm{TypeForAlt}$ field.
UPA17 therefore supplies precisely
$H.\mathrm{AutomorphicHeckeSet}
(\iota,\mathrm{UnitaryPointProperties}(\bar r,\iota))$.
AFS5, applied to its actual $\mathrm{IsTheta}$ field, proves
$\mathrm{DenseWt}$ for every weight.  The source admitted set
for this $hset$ is the entire set of weights, so the constant zero
dominant weight is an allowed initial $\lambda_0$.

Apply HEP11 to these original data and the given $D,\psi$.
This proves (HEP.24).  The four assertions about the point
representation are exactly the four fields of the supplied
$\mathrm{UnitaryPointProperties}$, already installed by UPA15--UPA17.
The determinant, local characters, level, and Hecke algebra remain
the ones in (HEP.1)--(HEP.2); there is no new coefficient algebra
or kernel ideal in this construction.  The statement makes no
claim about a level-to-level comparison, the complete
$\mathrm{RingProps}$ record, or an all-level automorphy endpoint.
Those are separate consumers of this full output.
\end{proof}
% END HEP12
